% Part: incompleteness
% Chapter: incompleteness-provability
% Section: introduction

\documentclass[../../../include/open-logic-section]{subfiles}

\begin{document}

\olfileid{inc}{inp}{int}

\olsection{परिचय}

हिल्बर्टच्या मते, नैसर्गिक संख्यांसारख्या गणितीय रचनेसाठीची
स्वयंसिद्धकपद्धती त्या रचनेविषयीची सर्व सत्य विधाने सिद्ध करू देत
नसेल, तर ती अपुरी आहे. निष्पत्तीच्या आकारिक पद्धतींविषयीच्या
त्याच्या नंतरच्या स्वारस्याचा विचार केला, तर नैसर्गिक संख्यांविषयी
युक्तिवाद करण्यासाठी वापरली जाणारी आकारिक पद्धत केवळ सुसंगतच
नव्हे, तर \emph{संपूर्ण}ही असावी, असे त्याला वाटत होते असे दिसते:
त्या भाषेतील प्रत्येक विधान किंवा त्याचा निषेध !!{derivable}
असावा. गोडेलचे पहिले अपूर्णता प्रमेय दाखवते की अशी
स्वयंसिद्धकपद्धती अस्तित्वात नाही: अंकगणितासाठी संपूर्ण, सुसंगत
आणि !!{axiomatizable} अशी कोणतीही आकारिक पद्धत नाही. खरे तर,
“पुरेशी प्रबळ”, सुसंगत आणि !!{axiomatizable} अशी कोणतीही
गणितीय उपपत्ती संपूर्ण नसते.

हिल्बर्टचा आणखी महत्त्वाचा उद्देश म्हणजे अभिजात गणिताला आधार
देण्याच्या कार्यक्रमाचा मुख्य भाग: अभिजात युक्तिवादाचे निरूपण
करणाऱ्या आकारिक पद्धतींच्या सुसंगततेच्या सांतवादी सिद्धता
शोधणे. म्हणून हिल्बर्टच्या कार्यक्रमासाठी गोडेलचे दुसरे
अपूर्णता प्रमेय हा अधिक मोठा धक्का होता. पहिल्या प्रमेयाप्रमाणे
दुसरेही प्रमेय काहीसे स्थूल शब्दांत सांगता येते. साधारणपणे ते
असे म्हणते की अंकगणिताची कोणतीही सुसंगत, प्रभावीपणे
स्वयंसिद्धकीकरण करता येणारी आणि पुरेशी प्रबळ उपपत्ती स्वतःची
सुसंगतता व्यक्त करणारे प्रमाण आकारिक वाक्य सिद्ध करू शकत नाही. येथे “पुरेशी प्रबळ” यामध्ये
केवळ~$\Th{Q}$ पेक्षा थोडे अधिक सामर्थ्य समाविष्ट करावे लागेल.

गोडेलच्या अपूर्णता प्रमेयाच्या मूळ सिद्धतेमागची कल्पना
एपिमेनिडीजच्या विरोधाभासात दिसते. क्रीटचा रहिवासी असलेल्या
एपिमेनिडीजने सर्व क्रीटवासी खोटारडे असतात असे म्हटले;
“हे वाक्य असत्य आहे” हे विधान त्या विरोधाभासाचे अधिक थेट
रूप आहे. सारतः, सत्याच्या जागी !!{derivability} ठेवून गोडेलने
स्वतःच !!{derivable} नसल्याचे अप्रत्यक्षपणे सांगणारे
!!a{sentence} आकारिक रीतीने तयार केले. ते !!{sentence}
!!{derivable} असेल, तर उपपत्ती विसंगत ठरेल. गोडेलने
विचारात घेतलेल्या स्वयंसिद्धकपद्धतीत त्या !!{sentence} चा
निषेधही !!{derivable} नसतो, असे त्याने दाखवले. (या दुसऱ्या
भागासाठी त्याला उपपत्ती~$\Th{T}$ ही तथाकथित
“$\omega$-सुसंगत” आहे असे गृहीत धरावे लागले.
$\omega$-सुसंगतता ही सुसंगततेशी संबंधित पण अधिक प्रबळ अट
आहे.\footnote{म्हणजे प्रत्येक $\omega$-सुसंगत उपपत्ती
सुसंगत असते; उलट विधान खरे नसते.} गोडेलनंतर काही वर्षांनी
रोसरने केवळ~$\Th{T}$ च्या साध्या सुसंगततेचे गृहीतक पुरेसे
आहे असे दाखवले.)

पहिली अडचण म्हणजे स्वतःचा निर्देश करणारे !!a{sentence}
कसे तयार करायचे हे समजणे. $\Th{Q}$ च्या भाषेतील प्रत्येक
!!{formula}~$!A$ साठी~$\gn{!A}$ हा~$\Gn{!A}$ या संख्येला
अनुरूप संख्यांक मानू. याचा अर्थ नीट पाहू: $!A$ हे
$\Th{Q}$ च्या भाषेतील !!a{formula} आहे, $\Gn{!A}$ ही
नैसर्गिक संख्या आहे, आणि $\gn{!A}$ हे $\Th{Q}$ च्या
भाषेतील \emph{पद} आहे. त्यामुळे $\Th{Q}$ च्या भाषेतील
प्रत्येक !!{formula}~$!A$ ला~$\gn{!A}$ असे \emph{नाव}
मिळते; ते $\Th{Q}$ च्या भाषेतील पद आहे. यामुळे
$\Th{Q}$ च्या भाषेतील !!{formula}s ना
इतर !!{formula}s विषयी काहीतरी “सांगता” येते अशी
संकल्पनात्मक चौकट मिळते. पुढील उपपत्तीला स्थिरबिंदू
उपपत्ती म्हणतात.

\begin{lem}
$\Th{T}$ ही~$\Th{Q}$ चा विस्तार असलेली कोणतीही उपपत्ती
असू द्या, आणि $!B(x)$ हे फक्त~$x$ हेच मुक्त चल असलेले
कोणतेही !!{formula} असू द्या. मग असे
!!a{sentence}~$!A$ असते की
$\Th{T} \Proves!A \liff !B(\gn{!A})$.
\end{lem}

ही उपपत्ती असे सांगते की कोणत्याही गुणधर्मासाठी $!B(x)$,
“$!B(x)$ माझ्याबाबतीत खरा आहे” असे सांगणारे
!!a{sentence}~$!A$ असते, आणि $\Th{T}$ ला हे
“माहीत” असते.

असे !!a{sentence} कसे तयार करायचे? क्वाइनने मांडलेले
एपिमेनिडीजच्या विरोधाभासाचे पुढील रूप पाहू:
\begin{quote}
“उद्धृत रूप आधी जोडल्यास असत्य विधान मिळते” उद्धृत रूप
आधी जोडल्यास असत्य विधान मिळते.
\end{quote}
हे वाक्य थेट स्वतःचा निर्देश करत नाही. ते फक्त अवतरणचिन्हांतील
विन्यासात्मक वस्तूंविषयी विधान करते; त्या अर्थाने ते
पुढील वाक्यांसारखेच आहे:
\begin{enumerate}
\item “रॉबर्ट” हे छान नाव आहे.
\item “मी धावलो.” हे छोटे वाक्य आहे.
\item “तीन शब्द आहेत” यात तीन शब्द आहेत.
\end{enumerate}
पण “उद्धृत रूप आधी जोडल्यास असत्य विधान मिळते” या
वाक्यांशाच्या आधी त्याचेच उद्धृत रूप जोडले, तर काय होते?
मूळ वाक्यच मिळते!{} थोडक्यात, ते वाक्य स्वतःच असत्य
असल्याचे सांगते.

\end{document}
